% Texto base para apoiar a escrita do TCC
% Fonte: analise dos notebooks do diretorio Laboratorio/

\section{Visao Geral do Laboratorio Experimental}

O laboratorio foi estruturado para comparar duas abordagens de deteccao de intrusao em rede: uma abordagem de Deep Learning (DNN) e uma abordagem de Machine Learning classica (Random Forest), ambas treinadas sobre o dataset CICIDS2017 apos etapa de tratamento de dados. O fluxo experimental foi dividido em tres cenarios: (i) cenario completo, com todas as classes de ataque presentes no treinamento; (ii) cenario sem PortScan no treinamento; e (iii) cenario sem XSS no treinamento. Essa organizacao permite avaliar tanto desempenho em cenario fechado (closed-set) quanto comportamento com ataques ausentes no treino.

% [INSERIR FIGURA 1: Fluxograma geral do pipeline - tratamento, treino e avaliacao]

\section{Objetivo de Cada Notebook e Papel no Experimento}

\subsection{Tratamento de Dados.ipynb}

Este notebook tem como objetivo consolidar e preparar os dados brutos do CICIDS2017 para os experimentos de classificacao. A etapa contempla leitura dos arquivos CSV, unificacao, remocao de inconsistencias e geracao dos conjuntos finais de treino e teste para cada cenario.

Resultados observados no output salvo do notebook:
\begin{itemize}
\item Tamanho total do dataset unificado: (2.830.743, 79).
\item Cenario completo: treino com 1.979.513 amostras e teste com 848.363 amostras.
\item Cenario sem PortScan: 110.968 registros PortScan movidos do treino para o teste; treino final com 1.868.545 amostras e teste com 959.331 amostras.
\item Arquivos finais salvos em \texttt{dataset tratado/} para alimentar os notebooks de treinamento.
\end{itemize}

% [INSERIR FIGURA 2: Tabela com quantidade de amostras por cenario]
% [INSERIR FIGURA 3: Distribuicao de classes antes e depois da separacao dos cenarios]

\subsection{Treinamento DNN.ipynb}

Este notebook implementa a rede neural densa para classificacao multiclasses no cenario completo. O objetivo principal e obter um baseline de desempenho da DNN quando todas as classes estao presentes no treinamento.

Etapas centrais: leitura dos dados tratados, codificacao dos rotulos, definicao da arquitetura da rede, treinamento com validacao e avaliacao por matriz de confusao e relatorio de metricas.

Resultado observado no output salvo:
\begin{itemize}
\item Acuracia global reportada: 0,98 (suporte de 848.363 amostras no teste).
\end{itemize}

% [INSERIR FIGURA 4: Arquitetura da DNN (camadas, ativacoes e dropout)]
% [INSERIR FIGURA 5: Matriz de confusao da DNN no cenario completo]

\subsection{Treinamento DNN - Sem Portscan.ipynb}

Este notebook avalia a DNN em cenario com classe ausente no treinamento (PortScan), com objetivo de medir degradacao de desempenho e comportamento de generalizacao quando aparece um ataque nao visto durante o aprendizado.

Resultado observado no output salvo:
\begin{itemize}
\item Acuracia global reportada: 0,79 (suporte de 959.331 amostras no teste).
\end{itemize}

Interpretacao: houve queda relevante de desempenho em relacao ao cenario completo, o que e coerente com o aumento de dificuldade quando uma classe de ataque e removida do treinamento.

% [INSERIR FIGURA 6: Matriz de confusao da DNN sem PortScan no treino]
% [INSERIR FIGURA 7: Comparacao de acuracia DNN - completo vs sem PortScan]

\subsection{Treinamento DNN - Sem XSS.ipynb}

Este notebook replica a avaliacao de generalizacao para outro ataque removido do treinamento (XSS), com a mesma estrutura de pipeline da DNN.

Resultado observado no output salvo:
\begin{itemize}
\item Acuracia global reportada: 0,98 (suporte de 848.822 amostras no teste).
\end{itemize}

Interpretacao: neste cenario, a queda de desempenho nao foi acentuada como no caso PortScan, sugerindo que o comportamento de generalizacao depende da classe removida e da proximidade estatistica entre classes.

% [INSERIR FIGURA 8: Matriz de confusao da DNN sem XSS no treino]

\subsection{Treinamento Random Forest.ipynb}

Este notebook treina o classificador Random Forest no cenario completo para estabelecer comparacao direta com a DNN, com foco em desempenho de classificacao e robustez do metodo.

Resultado observado no output salvo:
\begin{itemize}
\item Acuracia global reportada: 1,00 (suporte de 848.363 amostras no teste).
\end{itemize}

% [INSERIR FIGURA 9: Matriz de confusao do Random Forest no cenario completo]

\subsection{Treinamento Random Forest - Sem Portscan.ipynb}

Este notebook testa o Random Forest no cenario em que PortScan e removido do treinamento, buscando comparar a sensibilidade do modelo classico a ataques nao vistos.

Resultado observado no output salvo:
\begin{itemize}
\item Acuracia global reportada: 0,83 (suporte de 959.331 amostras no teste).
\end{itemize}

Interpretacao: ha queda de desempenho frente ao cenario completo, embora superior ao valor observado para a DNN no mesmo cenario.

% [INSERIR FIGURA 10: Matriz de confusao do Random Forest sem PortScan no treino]

\subsection{Treinamento Random Forest - Sem XSS.ipynb}

Este notebook aplica o mesmo protocolo para o cenario sem XSS no treinamento.

Resultado observado no output salvo:
\begin{itemize}
\item Acuracia global reportada: 1,00 (suporte de 848.822 amostras no teste).
\end{itemize}

Interpretacao: o resultado se manteve elevado neste cenario, indicando que a remocao de XSS do treino nao impactou fortemente a acuracia global.

% [INSERIR FIGURA 11: Matriz de confusao do Random Forest sem XSS no treino]

\section{Leitura Critica dos Resultados Observados}

Os resultados mostram que os modelos apresentam desempenho muito alto no cenario completo, com destaque para o Random Forest. Em cenarios com classe removida, a queda mais marcante ocorre no experimento sem PortScan, sugerindo maior sensibilidade do processo de classificacao a esse tipo de ataque quando ausente no treinamento. No experimento sem XSS, os dois modelos mantiveram desempenho global elevado.

Como limitacao metodologica, a acuracia global isolada nao descreve completamente o comportamento por classe. Por isso, na redacao final, recomenda-se enfatizar matriz de confusao e metricas por classe (precisao, recall e F1-score) para sustentar comparacoes de forma mais solida.

% [INSERIR FIGURA 12: Grafico comparativo de acuracia por modelo e cenario]
% [INSERIR FIGURA 13: Tabela de tempo de treino e inferencia por modelo]

\section{Esqueleto dos Capitulos que Devem Vir Antes do Laboratorio}

\subsection{Capitulo 1 - Fundamentacao do Problema}
\begin{itemize}
\item Contexto de seguranca cibernetica e crescimento de ataques em rede.
\item Papel de NIDS em ambientes corporativos e academicos.
\item Limites de deteccao baseada apenas em assinatura.
\item Motivacao para uso de IA na deteccao de intrusao.
\end{itemize}

\subsection{Capitulo 2 - Base Teorica para os Modelos}
\begin{itemize}
\item Conceitos de Machine Learning aplicados a classificacao de trafego.
\item Conceitos de Deep Learning para deteccao de padroes complexos.
\item Funcionamento do Random Forest e justificativa de escolha.
\item Funcionamento de redes neurais densas e justificativa de escolha.
\item Principais metricas: acuracia, precisao, recall, F1-score e matriz de confusao.
\end{itemize}

\subsection{Capitulo 3 - Base de Dados e Estrategia Experimental}
\begin{itemize}
\item Apresentacao do CICIDS2017 e tipos de ataque.
\item Descricao das features e desafios de dados (desbalanceamento e ruido).
\item Estrategia de tratamento dos dados (limpeza, normalizacao e split).
\item Definicao dos tres cenarios experimentais: completo, sem PortScan e sem XSS.
\item Hipoteses experimentais que serao testadas no laboratorio.
\end{itemize}

\subsection{Capitulo 4 - Metodologia de Implementacao}
\begin{itemize}
\item Ambiente computacional (bibliotecas e ferramentas).
\item Pipeline de treinamento e avaliacao para DNN e Random Forest.
\item Criterios de comparacao entre modelos.
\item Criterios de interpretacao dos resultados por cenario.
\end{itemize}

\section{Observacao Importante para a Redacao Final}

Os notebooks no estado atual aparecem como \textit{nao executados} no resumo do VS Code, mas possuem outputs salvos em celulas com resultados. Portanto, os valores aqui descritos foram extraidos desses outputs persistidos. Antes da versao final do texto, recomenda-se executar novamente todos os notebooks para congelar resultados atualizados no mesmo ambiente de execucao e reduzir variacoes de reproducao.
